\chapter{Protonation states and orientation}
\label{ch:protonation}

Before a simulation, you must decide which acidic and basic groups of the
protein carry a proton, and therefore a charge. This chapter explains in plain
words what that means and why the pH matters, then shows how ProPrep predicts
the states with PROPKA and writes the right residue names into the structure.
The last part covers the Structure Orientator, which turns the finished
structure to line it up with the axes.

\begin{glance}
\glanceitem{What it does}{Predicts the protonation state of every titratable
  residue at your pH, lets you accept or change each one, and writes a PDB file
  with the matching Amber residue names. Optionally rotates and centers the
  structure.}
\glanceitem{You need it when}{Always for a standard simulation: every residue
  needs a fixed charge. Orientation only when your setup needs a particular
  direction (membranes, fields, pulling).}
\glanceitem{Before you start}{Filter and repair first
  (Chapter~\ref{ch:cleaning}); the analyzer uses the filtered or repaired file
  automatically. Orient last.}
\glanceitem{Where to find it}{\menupath{4e}{Protonation State Analyzer} and
  \menupath{4f}{Structure Orientator}; ``fix'' stage in Workflow Mode.}
\glanceitem{You get}{\file{structure_with_md_names.pdb}, a PROPKA report in
  \file{propka_results/}, optional \file{ph_profile.csv}; from 4f a file ending
  in \file{_oriented.pdb}.}
\end{glance}

\section{What protonation means, and why pH matters}
\label{sec:protonation-why}

Some amino-acid side chains can take up or give away a proton (a hydrogen ion,
H$^+$). Gaining a proton adds a positive charge; losing one adds a negative
charge. Whether a group carries its proton is its \term{protonation state}. The
groups that can change are called \term{titratable} residues: aspartate (Asp),
glutamate (Glu), histidine (His), cysteine (Cys), tyrosine (Tyr), lysine (Lys),
arginine (Arg), and the two chain ends.

Each group has a \term{pKa}: the pH at which half of the molecules carry the
proton. When the pH is below the pKa (more acidic), the group mostly
\emph{has} its proton. When the pH is above the pKa, it mostly does not. The
textbook values ProPrep uses are:

\begin{tabularx}{\linewidth}{@{}lrL@{}}
\toprule
\textbf{Group} & \textbf{pKa} & \textbf{At pH 7, usually\ldots} \\
\midrule
Asp & 3.65 & without proton, charge $-1$ \\
Glu & 4.25 & without proton, charge $-1$ \\
His & 6.00 & close to its pKa: mostly neutral, partly charged \\
Cys & 8.18 & with proton, neutral \\
Tyr & 10.07 & with proton, neutral \\
Lys & 10.53 & with proton, charge $+1$ \\
Arg & 12.48 & with proton, charge $+1$ \\
\bottomrule
\end{tabularx}

\medskip
\noindent So the pH you choose decides the charges. At pH 4.5 more Asp and Glu
carry their proton and fewer are negative than at pH 7.

\textbf{Why this matters for a simulation.} A standard molecular-dynamics
simulation cannot add or remove protons while it runs. You fix every residue's
state at the start, and it stays that way. If you get an important one wrong,
for example in an active site, the simulation models a different chemistry.

\textbf{The protein changes the pKa.} The textbook values hold for a residue
in open water. Inside a protein, a residue may be buried away from water
(\term{desolvation}, which favors the uncharged form), or sit next to other
charges and hydrogen bonds. These can shift its pKa by several units.
\term{PROPKA} is a program that estimates these shifts from the structure. A
famous example is HIV protease: its two catalytic aspartates face each other,
and PROPKA predicts that one of them keeps its proton at pH 7 even though
aspartate normally loses it.

\textbf{Histidine has three forms.} Neutral histidine can carry its single
proton on either of two ring nitrogens, called ND1 (delta) and NE2 (epsilon).
These two forms are \term{tautomers}. With both nitrogens protonated, histidine
is positive. Amber force fields give each form its own residue name:
\term{HID} (proton on ND1), \term{HIE} (proton on NE2) and \term{HIP} (both,
charge $+1$). Which neutral form is right depends on the hydrogen bonds around
the ring; HIE is the more common.

\textbf{The probability threshold.} PROPKA gives each residue a
\emph{probability} of being protonated at your pH, between 0 and 1. ProPrep
must turn that into yes or no. With the default threshold of 0.5 it takes the
more likely state. A higher threshold (for example 0.7) assigns ``protonated''
only when it is quite sure; a lower one does so more readily. For residues near
0.5 neither answer is really right; for them, ProPrep's own advice is to
consider a constant-pH simulation (Section~\ref{sec:protonation-cph}).

\section{Step by step: predicting protonation states with PROPKA}
\label{sec:protonation-propka}

The example is the HIV protease dimer of \pdb{1HSG}, filtered to the two
protein chains as in Section~\ref{sec:cleaning-filter}.

\begin{steps}
  \item Type \key{4e}, then \key{1} (\prompt{Analyze protonation states
        (PROPKA)}). The other options stay blocked until this is done.
  \item Check the first line: \prompt{Using filtered structure:
        filtered\_structure.pdb}. The analyzer picks the most prepared file by
        itself (repaired before filtered before the original download).
  \item Type the pH of your experiment, or press \key{Enter} for 7.0. Then
        accept the threshold of 0.5 with \key{Enter}, and answer \key{y} to
        use PROPKA:
\begin{screen}
Using filtered structure: filtered_structure.pdb
No residues excluded from protonation analysis
Enter pH value for analysis (7.0):
...
Enter protonation probability threshold (0-1) (0.5):
Run PROPKA for improved pKa prediction? [y/n] (y): §ans«y»
\end{screen}
        PROPKA runs for a few seconds.
  \item Read the table. It is shown here with narrower columns and fewer rows:
\begin{screen}
          Protonation States at pH 7.0 (PROPKA pKa)
Index Chain    Residue    pKa   ΔpKa Desolv Protonated Probability Charge
    1     A   N_TERM 1   7.74  -0.26   0.43       True      0.8460   0.85
    2     A      ARG 8  14.00  +1.52   0.41       True      1.0000   1.00
...
    6     A     ASP 25   5.07  +1.42   2.99      False      0.0116  -0.99
...
   20     A     HIS 69   6.65  +0.65   0.36      False      0.3088   0.31
...
   23     A     CYS 95  11.51  +3.33   2.61       True      1.0000  -0.00
...
   30     B     ASP 25   9.06  +5.41   3.00       True      0.9914  -0.01
...
   44     B     HIS 69   6.02  +0.02   0.40      False      0.0948   0.09
...
   48     B  C_TERM 99   3.10      —      —      False      0.0001  -1.00
Net Charge: 5.13
\end{screen}
\end{steps}

\subsection{Reading the table}

\begin{description}
  \item[Index] The row number. You use it to ask for details (next step).
  \item[Chain, Residue] Which residue. \prompt{N\_TERM} and \prompt{C\_TERM} are
        the two ends of the chain.
  \item[pKa] PROPKA's estimate for this residue in this structure.
  \item[ΔpKa] The shift from the textbook value. A large shift means the
        surroundings matter a lot. Here ASP B25 is shifted by $+5.41$.
  \item[Desolv] The part of the shift caused by burial. Large positive values
        mean the residue is buried and prefers to be uncharged.
  \item[Protonated] ProPrep's decision at your pH and threshold.
  \item[Probability] The fraction of the time the residue carries the proton.
        Values far from 0 and 1 (HIS A69: 0.31) mark residues whose state is
        uncertain.
  \item[Charge] The average charge that follows from the probability.
  \item[Net Charge] The sum for the whole structure, here $+5.13$. The
        simulation will need counter-ions to balance it
        (Chapter~\ref{ch:topology}).
\end{description}

In \pdb{1HSG} the two catalytic aspartates differ: ASP A25 has a pKa of 5.07
and loses its proton, while ASP B25 has 9.06 and keeps it. This is the
``one protonated, one not'' pattern expected for this enzyme.

\subsection{Why does a residue have this pKa?}

After the table ProPrep asks
\prompt{Enter indices to visualize (e.g., 1,3,5-7), or Enter to skip}. Type an
index to see what drives that residue's pKa:
\begin{screen}
Enter indices to visualize (e.g., 1,3,5-7), or Enter to skip (): §ans«6»

ASP A:25  (predicted pKa 5.07, Δ +1.42 vs textbook)
  Buried 100 %; desolvation +2.99 (raises pKa)
┏━━━━━━━━━━━━━━┳━━━━━━━━━━━┳━━━━━━━━━━┳━━━━━━━━━━━━━━━┓
┃ Contribution ┃ Type      ┃ Partner  ┃ Effect on pKa ┃
┡━━━━━━━━━━━━━━╇━━━━━━━━━━━╇━━━━━━━━━━╇━━━━━━━━━━━━━━━┩
│        -0.89 │ sidechain │ ASP 25 B │ lowers (↓)    │
│        -0.60 │ backbone  │ GLY 27 A │ lowers (↓)    │
│        -0.17 │ backbone  │ ALA 28 A │ lowers (↓)    │
...
│        +0.03 │ coulombic │ ASP 29 A │ raises (↑)    │
└──────────────┴───────────┴──────────┴───────────────┘
✓ Structure viewer launched at http://localhost:8765/viewer (1 structure)
\end{screen}
Each row is one influence: a hydrogen bond from a side chain or the backbone,
or an electrostatic (\prompt{coulombic}) interaction with a charged neighbor.
Here ASP A25 is completely buried, which raises its pKa, but its hydrogen bond
to ASP B25 lowers it again. The 3D viewer opens with the residue in magenta,
partners that lower the pKa in green and those that raise it in orange.

\subsection{The titration curve}

Next, ProPrep offers a pH titration profile
(\prompt{Would you like to generate a pH titration profile? [y/n] (n)}). With
\key{y} and the default range (pH 0 to 14 in steps of 0.5) it prints the net
charge at each pH, a text plot, and the \term{isoelectric point} (pI, the pH at
which the net charge is zero): 9.37 for the protease dimer. You can then look up
values: type \typed{5.5} to get the charge at pH 5.5 (7.95), or \typed{0e} to
get the pH at charge 0. Type \typed{done} to finish. The profile can be saved
as \file{ph_profile.csv}.

The raw PROPKA output is kept in the folder \file{propka_results/} (files
ending in \file{.pka} and \file{.pdb}).

\section{Step by step: setting the residue names for the simulation}
\label{sec:protonation-names}

The force field recognizes each protonation state by a residue name. This step
writes those names into a new PDB file. Only the names change; no hydrogen
atoms are added yet (tLEaP adds them in Chapter~\ref{ch:topology}).

\begin{steps}
  \item In the Protonation State Analyzer menu, type \key{3}
        (\prompt{Set MD residue names}).
  \item At \prompt{Is this for constant pH simulation? [y/n] (n)}, press
        \key{Enter} (no) for a normal simulation.
  \item At \prompt{Set all residues to standard protonation states \ldots ?
        [y/n] (n)}, press \key{Enter} (no) to use PROPKA's predictions.
  \item ProPrep now goes through the residue types one by one: ASP, CYS, GLU,
        HIS, LYS, TYR. For each, it shows a table with a recommended name per
        residue and asks whether to use the recommendations:
\begin{screen}
                             ASP Residues
┏━━━━━━━┳━━━━━━━┳━━━━━━┳━━━━━━━━━━━━━━━━━┳━━━━━━━━━━━━━┳━━━━━━━━━━━━━┓
┃ Chain ┃ ResID ┃ pKa  ┃ State at pH 7.0 ┃ Probability ┃ Recommended ┃
┡━━━━━━━╇━━━━━━━╇━━━━━━╇━━━━━━━━━━━━━━━━━╇━━━━━━━━━━━━━╇━━━━━━━━━━━━━┩
│ A     │ 25    │ 5.07 │ Deprotonated    │ 0.012       │ ASP         │
...
│ B     │ 25    │ 9.06 │ Protonated      │ 0.991       │ ASH         │
...
└───────┴───────┴──────┴─────────────────┴─────────────┴─────────────┘

Options for ASP residues:
• ASP - Deprotonated carboxylate (COO⁻, -1 charge)
• ASH - Protonated carboxylic acid (COOH, neutral)
Note: ASP is typically deprotonated at physiological pH

Use recommended names for all ASP residues? (y):
\end{screen}
        Press \key{Enter} to accept. This is the right answer in most cases.
  \item At the end, ProPrep sums up and writes the file:
\begin{screen}
Set MD names for 44 residues

Summary of assignments:
  ARG: 8 residues
  ASH: 1 residues
  ASP: 7 residues
...
Update structure with these residue names? [y/n] (y): §ans«y»
Using analyzed structure for MD residue name update: filtered_structure.pdb
Enter output filename for updated structure (structure_with_md_names.pdb):

Modified 44 residue names in PDB file
...
Saved updated structure to structure_with_md_names.pdb
\end{screen}
\end{steps}

The file \file{structure_with_md_names.pdb} is the input for the next stages.
Arginine is never asked about; it is always ARG.

\subsection{The residue names}

\begin{tabularx}{\linewidth}{@{}llL@{}}
\toprule
\textbf{Residue} & \textbf{Name} & \textbf{Meaning} \\
\midrule
Asp & ASP / ASH & charged ($-1$) / with proton, neutral \\
Glu & GLU / GLH & charged ($-1$) / with proton, neutral \\
His & HIE / HID / HIP & proton on NE2 / proton on ND1 / both, charged ($+1$) \\
Cys & CYS / CYM & with proton (SH), neutral / without, charged ($-1$) \\
Lys & LYS / LYN & charged ($+1$) / without proton, neutral \\
Tyr & TYR & with proton, neutral \\
Arg & ARG & always charged ($+1$) \\
\bottomrule
\end{tabularx}

\medskip
Appendix~\ref{app:residues} lists all residue names used by Amber.

\subsection{Changing the recommendation for some residues}
\label{sec:protonation-override}

You may know better than PROPKA, for example from experiments on the active
site. Answer \key{n} at \prompt{Use recommended names for all \ldots{}
residues?}. ProPrep then explains the syntax:
\begin{screen}
Syntax: Enter a name to apply to all, or specify overrides:
  • Single name: ASH (applies to all)
  • Overrides: 94,119:ASP 122:ASH (applies to listed resids in every chain)
  • Chain-specific: A:94:ASP (restrict to one chain)
  • Press Enter to use all recommendations
Customize ASP names (): §ans«A:25:ASH B:25:ASP»
  A:25 → ASH
  B:25 → ASP
\end{screen}
Residues you do not name keep their recommendation. In this test we swapped the
protonated aspartate from chain B to chain A. For histidines, the same way, we
typed \typed{A:69:HIP B:69:HID} to make one histidine charged and the other a
delta tautomer. The file then contained ASH A25, HIP A69 and HID B69.

\begin{watchout}
An override that ProPrep cannot read is skipped, and ProPrep moves on to the
next residue type \emph{without asking again}:
\begin{screen}
Customize CYS names (): §ans«XYZ»
Skipping invalid syntax: XYZ (expected resid:NAME)
...
Customize GLU names (): §ans«21:LYS»
Invalid name 'LYS' - valid choices: GLH, GLU
\end{screen}
The recommendations are then used for that type. If this happens, run option
\key{3} again. Residue numbers are those of the analyzed file; after the
Structure Fixer they start from 1 (Section~\ref{sec:cleaning-caps}).
\end{watchout}

\begin{watchout}
Two more things to know in version 1.23.0:
\begin{itemize}
  \item Running option \key{3} again overwrites
        \file{structure_with_md_names.pdb} without asking. Type a different
        file name if you want to keep both.
  \item After overrides, option \key{4} (View) shows the new names, but the
        Protonated, Probability and Charge columns and the net charge still
        show PROPKA's values. For example \prompt{ASH 25} appears with
        \prompt{False} and $-0.99$. The export (option \key{5}) does the same.
        The file written in this step has the right names.
\end{itemize}
\end{watchout}

\section{Constant-pH simulations}
\label{sec:protonation-cph}

In a \term{constant-pH simulation}, chosen residues are allowed to gain and
lose protons during the run, so their states follow the pH instead of being
fixed. This is useful when a residue's probability is near 0.5. Amber marks
such residues with special names: \term{AS4} for aspartate, \term{GL4} for
glutamate and HIP for histidine.

Answer \key{y} at \prompt{Is this for constant pH simulation?}. ProPrep leaves
out residues that cannot titrate in this mode (arginine and the chain ends) and
then, for each type, asks which residues to titrate:
\begin{screen}
Select ASP residues to TITRATE (will use AS4 name):
Enter indices (e.g., 1,3 or 1-3), 'all', or 'none':
Indices to titrate (): §ans«1»
  B:25 → AS4 (titratable)

Setting names for non-titratable ASP residues:
...
Use recommended names for all non-titratable ASP residues? (y):
\end{screen}
The other residues of that type get fixed names as before. The result goes to
the same file as before,
\file{structure_with_md_names.pdb}, and overwrites it. Running the simulation itself is covered in
Chapter~\ref{ch:runningmd}.

\begin{example}
Protonation depends on everything you kept. In this test the filter had kept
only chain B of the protease. Without its partner, ASP B25 got a pKa of 3.8 and
was \emph{deprotonated}, against 9.06 in the dimer. Filter first, and filter
carefully, before you protonate.
\end{example}

\section{Other choices}
\label{sec:protonation-other}

\paragraph{Textbook values instead of PROPKA.} Answer \key{n} at
\prompt{Run PROPKA for improved pKa prediction?}. The table then uses the
textbook pKa values (title \prompt{(Textbook pKa)}, without the ΔpKa and Desolv
columns). Tried on ubiquitin (\pdb{1UBQ}) at pH 4.5 with a threshold of 0.7.

\paragraph{Standard states for everything.} At \prompt{Set all residues to
standard protonation states (ASP/GLU deprotonated, CYS/LYS/ARG/TYR protonated,
HIS->HIE, heme PRN->PRD)?}, answer \key{y}. ProPrep ignores the predictions and
applies these textbook states. You can then type a file name, for example
\typed{ubq\_standard.pdb}. This is a quick choice for a protein at neutral pH
without unusual sites.

\paragraph{Charge capacitance (option 2).} Calculates how strongly the
protein's charge responds to small pH changes, per residue type and in total.
It can also draw this as a curve over pH and save it as
\file{capacitance_profile.csv}. Most users do not need it.

\paragraph{View and export (options 4 and 5).} Option \key{4} shows the table
again. Option \key{5} asks \prompt{Select output format (text)} but does not
list the choices. Type \typed{text} (the default) or \typed{csv}. The files are
called \file{protonation_analysis.text} and
\file{protonation_analysis.csv}.

\paragraph{Metal and redox sites.} If redox sites were detected earlier
(Chapter~\ref{ch:redox}), the analysis first lists them and asks which to leave
out of the analysis. Residues that bind a metal have strongly shifted pKa
values, so excluding them is often sensible.

\section{Orienting the structure}
\label{sec:protonation-orient}

The Structure Orientator (\key{4f}) rotates and moves the whole structure. The
chemistry does not change, only the coordinates. You need it when the direction
of the protein matters: for example a membrane protein that must stand upright
along the z axis (Chapter~\ref{ch:membrane}), a simulation with an electric
field, or pulling along one axis. For an ordinary protein in water you can skip
it.

\begin{watchout}
\textbf{Orient last.} The orientator always prefers its own earlier output: the
first line says \prompt{Using Oriented PDB file: \ldots}. The Protonation State
Analyzer, on the other hand, ignores oriented files. So if you orient and then
filter or protonate again, the orientator keeps offering the \emph{old}
oriented file, with the old chains and names. In our test, after a second
filter (chain B only) and a new protonation, \key{4f} still used the oriented
two-chain file. If you change anything earlier, start a new project or orient
again from the new file.
\end{watchout}

\begin{steps}
  \item Type \key{4f}, then \key{1}. Check the input file and choose a method:
\begin{screen}
Using Protonation-Updated PDB file: structure_with_md_names.pdb

Input structure: structure_with_md_names.pdb
Source: protonation_pdb_file
✓ Structure viewer launched at http://localhost:8765/viewer (1 structure)

Alignment options:
  1. Align principal axes to Cartesian axes (PCA)
  2. Align vector between two residues to an axis
  3. Align vector between two atoms to an axis
  4. Align plane (from residues) normal to Z
  5. Center structure at origin
  6. Apply custom rotation (Euler angles)
  7. Cancel

Select alignment method (1): §ans«1»
\end{screen}
  \item For method \key{1}, answer which axis the long axis of the protein
        should point along, then the middle axis. The shortest axis takes the
        remaining direction. Then choose whether to center:
\begin{screen}
Principal moments of inertia:
  Longest axis  (smallest moment): λ = 21322.2
  Middle axis:                      λ = 42699.6
  Shortest axis (largest moment):   λ = 50600.1

Map principal axes to Cartesian axes:
Align LONGEST axis to [X/Y/Z] (Z):
Align MIDDLE axis to [X/Y] (X): §ans«Y»
Shortest axis will align to X
Center structure at origin? [y/n] (y): §ans«y»

Applying transformation...
✓ Transformed 1,514 atoms
✓ Saved: structure_with_md_names_oriented.pdb
Transformation saved: structure_with_md_names_oriented.orientation.json
\end{screen}
\end{steps}

The \term{principal axes} are the protein's natural long, middle and short
directions, worked out from how its mass is spread. Aligning them with the x,
y and z axes makes the protein lie ``straight'' in the box. The
\file{.orientation.json} file records exactly what was done (method, rotation,
shift), so the step can be traced later. The viewer opens before and after,
with reference axes.

\subsection{The other methods}

\begin{tabularx}{\linewidth}{@{}lL@{}}
\toprule
\textbf{Method} & \textbf{What you give, and what it does} \\
\midrule
\key{1} Principal axes & Axis for the longest and the middle direction;
  center or not. \\
\key{2} Two residues & Chain and number of a first residue (placed at the
  negative end) and a second (positive end), then the axis. Example: A25 to
  B25, 6.04~Å apart, along X. \\
\key{3} Two atoms & As 2, plus an atom name (default CA; the second atom
  defaults to the first atom's name). Example: CG of A25 to CG of B25,
  5.17~Å, along Y. \\
\key{4} Plane to Z & At least three residues as \typed{A:10, A:20, B:10,
  B:20}; the plane through them is turned to face the z axis. Meant for
  membrane-facing residues. With fewer than three: \prompt{Need at least 3
  residues to define a plane.} \\
\key{5} Center only & Moves the center to the origin, by mass (\key{y}) or
  by geometry (\key{n}). No rotation. \\
\key{6} Euler angles & Rotations in degrees around Z, then Y, then X; option to
  center first. \\
\key{7} Cancel & Nothing is written. \\
\bottomrule
\end{tabularx}

\medskip
A residue or atom that does not exist gives \prompt{Residue B:500 not found.}
or \prompt{Atom B:25:XX not found.} and returns to the menu.

\begin{watchout}
Each run takes the previous oriented file as its input and adds another
\file{_oriented} to the name. After six runs in our test the name contained
\file{_oriented} six times (\file{structure_with_md_names_oriented_..._oriented.pdb}).
Rotations add up in the same way. If you want to try several orientations,
note that each one starts from the last, not from the original.
\end{watchout}

\section{If something goes wrong}
\label{sec:protonation-problems}

\begin{description}[style=nextline]
  \item[``Invalid pH value, using default of 7.0''] Anything that is not a
        number is replaced by 7.0 without asking again. The same happens to
        the threshold (\prompt{Invalid threshold, using default of 0.5}), for
        example if you type \typed{20}. If you meant another value, run option
        \key{1} again.
  \item[An override did not take effect] Invalid overrides are skipped without
        asking again (Section~\ref{sec:protonation-override}). Run option
        \key{3} again.
  \item[The View table does not show my override] Known limitation: the
        names are updated, the charges are not. The written PDB file is
        correct.
  \item[Residue numbers do not match the original] The Structure Fixer
        renumbered the chain from 1 (Section~\ref{sec:cleaning-caps}).
  \item[The analyzer ignores my latest structure] It uses repaired, then
        filtered, then downloaded files, and never oriented ones. Read the
        \prompt{Using \ldots} line.
  \item[The orientator uses an old file] It prefers any earlier oriented
        file. Orient last, or use a new project.
  \item[``Received 'Z', which is not one of the available options''] For the
        middle axis only the two axes not yet used are allowed.
\end{description}

See also Chapters~\ref{ch:troubleshooting} and~\ref{ch:knownissues}.
